\lecture{1}{2026-09-09}{Combinatorial constructions}

\begin{definition}[Combinatorial class]
A \emph{combinatorial class} is a pair \((\mathcal{C}, |\cdot|)\) where \(\mathcal{C}\) is a set and \(|\cdot|\colon \mathcal{C}\to \mathbb{N}\) is a \emph{weight function} such that for each \(n\in \mathbb{N}\), the set
\[
    \mathcal{C}_n = \{c\in \mathcal{C} : |c|=n\}
\]
is finite.
The sequence \(c_n = |\mathcal{C}_n|\) for \(n \ge 0\) is the \emph{counting sequence} of \(\mathcal{C}\).
\end{definition}

\begin{definition}[Generating function]
The \emph{generating function} of a combinatorial class \(\mathcal{C}\) is the formal power series
\[
    C(z) = \sum_{n\ge 0} c_n z^n = \sum_{c\in \mathcal{C}} z^{|c|}.
\]
\end{definition}

\begin{example}[Binary strings]
The set \(\mathcal{C} = \{\varepsilon, 0, 1, 00, 01, 10, 11, 000, \dots\}\) of all binary strings with the weight function \(|\cdot|\) being the length of the string is a combinatorial class with \(c_n = 2^n\). Its generating function is
\[
    C(z) = \sum_{n\ge 0} 2^n z^n = \frac{1}{1-2z}.
\]
However, the same set with weight function being the number of zeroes is not a combinatorial class, since \(\mathcal{C}_0 = \{\varepsilon, 1, 11, 111, \dots\}\) is infinite.
\end{example}

\section*{Combinatorial constructions}

To build generating functions without first calculating counting sequences, we define combinatorial operations on classes that translate directly into algebraic operations on their generating functions.

\subsection*{Base classes}
\begin{description}
    \item[Neutral class \(\mathcal{E}\):] \(\mathcal{E} = \{\varepsilon\}\) contains one object of size \(0\). Generating function: \(E(z) = 1\).
    \item[Atomic class \(\mathcal{Z}\):] \(\mathcal{Z} = \{\bullet\}\) contains one object of size \(1\). Generating function: \(Z(z) = z\).
\end{description}

\subsection*{Basic constructions}
Let \(\mathcal{A}\) and \(\mathcal{B}\) be combinatorial classes with generating functions \(A(z)\) and \(B(z)\).
\begin{description}
    \item[Sum (\(\mathcal{A} + \mathcal{B}\)):]
    The disjoint union of \(\mathcal{A}\) and \(\mathcal{B}\), with weight function \(|c| = |a|\) if \(c = a \in \mathcal{A}\) and \(|c| = |b|\) if \(c = b \in \mathcal{B}\).
    \[
        C(z) = A(z) + B(z).
    \]
    \item[Product (\(\mathcal{A} \times \mathcal{B}\)):]
    The Cartesian product of \(\mathcal{A}\) and \(\mathcal{B}\), with \(|(a,b)| = |a| + |b|\).
    \[
        C(z) = A(z) B(z).
    \]
    \item[Sequence (\(\operatorname{SEQ}(\mathcal{A})\)):]
    Assuming \(\mathcal{A}_0 = \varnothing\), \(\operatorname{SEQ}(\mathcal{A}) = \mathcal{E} + \mathcal{A} + \mathcal{A}^2 + \mathcal{A}^3 + \dots\) with \(|(a_1, \dots, a_k)| = \sum_{i=1}^k |a_i|\).
    \[
        C(z) = \sum_{k\ge 0} A(z)^k = \frac{1}{1 - A(z)}.
    \]
\end{description}

\section*{Iterative specifications}

\begin{example}[Binary strings]
Let \(\mathcal{Z}_0, \mathcal{Z}_1\) be atomic classes representing \(0\) and \(1\) respectively (so \(Z_0(z) = Z_1(z) = z\)).
The class of binary strings is
\[
    \mathcal{B} = \operatorname{SEQ}(\mathcal{Z}_0 + \mathcal{Z}_1).
\]
Its generating function is
\[
    B(z) = \frac{1}{1 - (Z_0(z) + Z_1(z))} = \frac{1}{1 - 2z}.
\]
\end{example}

\begin{example}[Compositions]
A \emph{composition} of \(n\) is a sequence of positive integers summing to \(n\).
Let \(\mathcal{P}\) be the class of positive integers (with \(|n| = n\)), and let \(\mathcal{C}\) be the class of compositions.
Then
\[
    \mathcal{P} = \mathcal{Z} \times \operatorname{SEQ}(\mathcal{Z}) \implies P(z) = \frac{z}{1-z},
\]
and
\[
    \mathcal{C} = \operatorname{SEQ}(\mathcal{P}) \implies C(z) = \frac{1}{1 - P(z)} = \frac{1}{1 - \frac{z}{1-z}} = \frac{1-z}{1-2z} = 1 + \frac{z}{1-2z}.
\]
Extracting the coefficient for \(n \ge 1\):
\[
    c_n = [z^n] C(z) = [z^n] \frac{z}{1-2z} = [z^{n-1}] \frac{1}{1-2z} = 2^{n-1}.
\]
\end{example}